% Copyright 2026 MarcosHCK
%
\documentclass{article}[12pt,a4paper]
\usepackage{amsfonts}
\usepackage{amsmath}
\usepackage{amsthm}
\usepackage{bookmark}
\usepackage{hyperref}
\usepackage{cleveref} % must load after hyperref
\usepackage{fontawesome5}
\usepackage{graphicx}
\usepackage{inputenc}[utf8]
\usepackage{indentfirst}
\usepackage{minted}
\usepackage{polyglossia}[spanish,english]
\usepackage{biblatex}[style=apa] % must load after polyglossia
\usepackage{csquotes} % must load after biblatex
\usepackage{setspace}
\usepackage{svg}
\usepackage{xcolor}[dvipsnames]

\setmainlanguage{spanish}
\setotherlanguage{english}

\setstretch{1.5}

\newcommand{\bashps}{\textasciitilde\#\,}

\definecolor{Aquamarine}{rgb}{0.498, 1, 0.83}
\definecolor{DarkAquamarine}{rgb}{0.18, 0.55, 0.34}
\definecolor{YellowOrange}{rgb}{1, 0.7, 0}
\newcommand\myshade{85}
\colorlet{mycitecolor}{YellowOrange}
\colorlet{mylinkcolor}{violet}
\colorlet{myurlcolor}{DarkAquamarine}

\hypersetup {
    colorlinks = true,
    citecolor  = mycitecolor!\myshade!black,
    linkcolor  = mylinkcolor!\myshade!black,
    urlcolor   = myurlcolor!\myshade!black,
  }

\begin{document}

\begin{titlepage}

  \centering
  {\huge\bfseries Reporte (semana 5 y 6)\par}
  {\Large Autor: \itshape Marcos Antonio Pérez Lorenzo\par}
\end{titlepage}

\pagebreak
% Copyright 2026 MarcosHCK
%
\section{Arquitecturas generales}

\subsection{Arduino UNO}
\label{sec:genarc:arduino}

Todos los programas escritos para ejecutarse en un \textbf{Arduino UNO} R3 siguen las siguientes características:

\begin{itemize}
  \item Lenguaje: \textbf{C++}.
  \item Dependencias: compilador de C/C++ para la arquitectura AVR (específicamente el ATmega328P)
  \item Flujo: obtención de datos, cálculos, salida del programa.
\end{itemize}

\subsection{C++}
\label{sec:genarc:cplusplus}

Todos los programas escritos en \textbf{C++} siguen las siguientes características

\begin{itemize}
  \item Lenguaje: \textbf{C++}.
  \item Dependencias: librería estándar.
  \item Flujo: obtención de datos, cálculos, salida del programa.
\end{itemize}

\subsection{Python}
\label{sec:genarc:python}

Todos los programas escritos en \textbf{Python} siguen las siguientes características

\begin{itemize}
  \item Lenguaje: \textbf{Python}.
  \item Dependencias: librería estándar.
  \item Flujo: obtención de datos, cálculos, salida del programa.
\end{itemize}
% Copyright 2026 MarcosHCK
%
\section{Instrucciones de uso}

\subsection{Arduino UNO}
\label{src:genericinstrc:arduino}

Para ejecutar un sketch de Arduino, compilarlo para la placa correspondiente,
cargar el firmware y seguir la interfaz de entrada y salida que describe su
sección específica.

\subsection{C++}
\label{src:genericinstrc:cplusplus}

Todos los programas escritos en \textbf{C++} se pueden ejecutar mediante un script de \textbf{Python} intermediario para su fácil uso, por lo tanto se ejecutan de la misma forma usado para programas escritos en \textbf{Python} (descritos en \ref{src:genericinstrc:python})

\subsection{Python}
\label{src:genericinstrc:python}

Todos los programas escritos en \textbf{Python} se pueden ejecutar desde una terminal con acceso al interprete. Por ejemplo, para ejecutar el programa "clase567": en una terminal cuyo directorio de trabajo es el directorio raíz del repositorio, ejecutar "\mintinline{bash}{python3 clase567}".

\section{Programas}

Todos los programas se encuentran alojados en este
\href{https://github.com/MarcosHCK/master-programming/}{repositorio \faGithubSquare}.
En esta entrega se presentan dos entrenadores de perceptrones y cuatro sketches
de inferencia para Arduino UNO. Los programas \texttt{clase010} y
\texttt{clase011} entrenan modelos a partir de tablas de verdad y archivos CSV,
respectivamente; los sketches de los programas 3 a 6 ejecutan modelos lógicos o
modelos aprendidos en una placa.

\subsection{Programa 1: \texttt{clase010}}

\subsubsection{Objetivos}

El programa \textbf{clase010} permite entrenar y consultar un perceptrón para
una función lógica booleana. Sus objetivos son:

\begin{enumerate}
  \item Generar automáticamente todas las combinaciones binarias para la
    cantidad de entradas seleccionada.
  \item Construir los resultados esperados de las funciones \texttt{and},
    \texttt{or}, \texttt{nand}, \texttt{nor} y \texttt{xor}.
  \item Ajustar un perceptrón de una neurona y presentar la probabilidad y la
    clasificación obtenidas para cada combinación.
  \item Permitir la exportación opcional de los errores de entrenamiento a
    JSON y de los pesos aprendidos a un encabezado C.
\end{enumerate}

\subsubsection{Tecnicismos}

Arquitectura basada en \ref{sec:genarc:python}. El punto de entrada está en
\path{clase010/__main__.py}; el modelo reutilizable está definido en
\path{clase010/Perceptron.py} y las funciones de activación en
\path{clase010/sigmoid.py}. No requiere dependencias externas a la biblioteca
estándar de Python.

La función \mintinline{python3}{binary_pairs} genera \(2^n\) vectores binarios
de longitud \(n\). El diccionario \mintinline{python3}{logics} aplica la
operación seleccionada a cada vector para producir los objetivos de
entrenamiento. El clasificador calcula una suma ponderada más un sesgo y aplica
la función sigmoide:

\[
  p = \sigma\!\left(b + \sum_{i=1}^{n} w_i x_i\right),
  \qquad
  \sigma(z) = \frac{1}{1 + e^{-z}}.
\]

El entrenamiento actualiza los pesos y el sesgo con descenso de gradiente y
pérdida cuadrática. La parada anticipada cuenta las épocas sin una mejora de
pérdida mayor que \texttt{min\_delta}; cuando transcurren las épocas indicadas
por \texttt{patience}, recupera los parámetros del mejor punto observado. En
esta interfaz se supervisa la pérdida de entrenamiento. El argumento
\texttt{--patience 0} desactiva la parada anticipada.

\subsubsection{Código relevante}

La generación del conjunto de entrenamiento y la invocación del clasificador
son independientes del número de entradas:

\begin{minted}{python3}
inputs = [[float (value) for value in row]
          for row in binary_pairs (args.inputs)]
targets = logics [args.logic] (inputs)

net = PerceptronClassifier (
  n_inputs = len (inputs [0]),
  learning_rate = args.learning_rate,
  activation = sigmoid,
  derivative = sigmoid_derivative_from_output)

errors = net.fit (
  inputs,
  targets,
  epochs = args.epochs,
  patience = None if 0 == args.patience else args.patience,
  min_delta = args.min_delta)
\end{minted}

El método \mintinline{python3}{fit} devuelve una lista con la pérdida media de
cada época ejecutada. La opción \mintinline{bash}{--errors-output} escribe esta
lista como un arreglo JSON. La opción \mintinline{bash}{-h} o
\mintinline{bash}{--header} guarda en un encabezado C el arreglo
\mintinline{cpp}{weights}, con el sesgo en la primera posición y luego los
pesos de las entradas.

\subsubsection{Ejecución}

Ver sección \ref{src:genericinstrc:python}. Desde la raíz del repositorio, por
ejemplo:

\begin{minted}{bash}
python -m clase010 --logic and --inputs 2 --epochs 5000 \\
  --patience 10 --errors-output errores.json \\
  --header modelo.h
\end{minted}

La ayuda del programa presenta la lista completa de opciones y sus descripciones.

\subsubsection{Resultados esperados}

La ejecución presenta la pérdida durante el entrenamiento, una tabla con las
entradas, el objetivo, la probabilidad predicha y la clase binaria, y finalmente
el sesgo y los pesos. Un ejemplo de la tabla para la función AND tiene esta
estructura:

\begin{verbatim}
 x1 x2 | target |   prob | pred
 ------+--------+--------+------
  0  0 |      0 | 0.08   |    0
  0  1 |      0 | 0.29   |    0
  1  0 |      0 | 0.29   |    0
  1  1 |      1 | 0.65   |    1
\end{verbatim}

Los valores numéricos dependen de la inicialización aleatoria y de las opciones
de entrenamiento. Una neurona lineal con sigmoide puede representar AND, OR,
NAND y NOR; XOR no es linealmente separable, por lo que el clasificador de una
neurona no puede ajustarse perfectamente a esa tabla.

\subsubsection{Liga de \textbf{GitHub}}

\href{https://github.com/MarcosHCK/master-programming/tree/master/clase010}{Aquí \faGithubSquare}

\subsection{Programa 2: \texttt{clase011}}

\subsubsection{Objetivos}

El programa \textbf{clase011} aplica el mismo clasificador a datos externos,
permitiendo que cada ejemplo tenga cualquier cantidad positiva de entradas
numéricas. Sus objetivos son:

\begin{enumerate}
  \item Leer filas de un CSV donde las primeras \(n-1\) columnas representan
    las entradas y la última columna contiene la clase.
  \item Detectar las dos etiquetas presentes en los datos y asociarlas de forma
    determinista con los objetivos \(0\) y \(1\).
  \item Normalizar las entradas por columna, entrenar el perceptrón y mostrar
    las predicciones y la exactitud sobre los ejemplos suministrados.
  \item Permitir exportar los errores a JSON y los parámetros del modelo y de
    normalización a un encabezado C.
\end{enumerate}

\subsubsection{Tecnicismos}

Arquitectura basada en \ref{sec:genarc:python}. El cargador usa el módulo
\mintinline{python3}{csv} de la biblioteca estándar. Cada fila debe tener el
mismo ancho y al menos dos columnas: una o más entradas finitas que puedan
convertirse a números, y una etiqueta de texto no vacía. Se requiere que el
conjunto contenga exactamente dos etiquetas distintas; el programa las ordena
alfabéticamente y muestra cuál se codifica como cero y cuál como uno.

Cada entrada \(x_i\) se centra y escala usando la media \(\mu_i\) y la
desviación estándar poblacional \(s_i\) calculadas sobre los ejemplos:

\[
  x_i' = \frac{x_i - \mu_i}{s_i}.
\]

Si una columna es constante, se usa escala uno para evitar la división entre
cero. El número de pesos se determina con el ancho leído del CSV. Al exportar
el modelo se escriben \mintinline{cpp}{weights} (sesgo primero),
\mintinline{cpp}{means} y \mintinline{cpp}{scales}; para evaluar nuevas entradas
se deben normalizar con las medias y escalas exportadas antes de aplicar la
suma ponderada y la sigmoide. En la interfaz de consola, la pérdida observada
para la parada anticipada es la pérdida de entrenamiento.

\subsubsection{Código relevante}

El lector separa la última columna como etiqueta y convierte todas las columnas
anteriores en el vector de entrada:

\begin{minted}{python3}
inputs = [float (field) for field in fields [:-1]]
label = fields [-1]

rows.append ((line, inputs, label))
\end{minted}

Una vez cargados los ejemplos, la cantidad de entradas se obtiene de la primera
fila y se usa para inicializar el modelo:

\begin{minted}{python3}
inputs = [row [1] for row in rows]
targets = [float (labels.index (row [2])) for row in rows]
normalized, means, scales = normalize_inputs (inputs)

net = PerceptronClassifier (
  n_inputs = len (inputs [0]),
  learning_rate = args.learning_rate,
  seed = args.seed,
  activation = sigmoid,
  derivative = sigmoid_derivative_from_output)
\end{minted}

La opción \mintinline{bash}{--csv-header} controla la fila de encabezado del
CSV (detección automática por omisión, o los valores \texttt{yes} y
\texttt{no}). La opción \mintinline{bash}{-h} o
\mintinline{bash}{--header} es distinta: recibe la ruta de salida del encabezado
C. El archivo exportado contiene las ponderaciones y también los datos
necesarios para reproducir la normalización.

\subsubsection{Ejecución}

Ver sección \ref{src:genericinstrc:python}. Un ejemplo de archivo de entrada
con dos entradas y una etiqueta es:

\begin{minted}{text}
edad,puntaje,grupo
18,20,joven
20,25,joven
60,80,adulto
65,90,adulto
\end{minted}

Desde la raíz del repositorio se puede entrenar y exportar los resultados con:

\begin{minted}{bash}
python -m clase011 clases.csv --epochs 5000 --patience 10 \\
  --errors-output errores.json --header modelo.h
\end{minted}

La opción \mintinline{bash}{--help} muestra el resto de las opciones,
como \mintinline{bash}{--csv-header}, que selecciona el modo de lectura del
encabezado CSV, y \mintinline{bash}{--seed} y \mintinline{bash}{--min-delta},
que controlan la semilla y la mejora mínima de pérdida, respectivamente.

\subsubsection{Resultados esperados}

El programa informa la correspondencia entre etiquetas y objetivos, muestra la
pérdida de entrenamiento y presenta una tabla con cada vector, la etiqueta
real, la probabilidad, la etiqueta predicha y si la predicción fue correcta.
En el ejemplo, la salida identifica \texttt{adulto} como clase cero y
\texttt{joven} como clase uno; la exactitud final se calcula sobre las mismas
filas utilizadas para entrenar, no sobre un conjunto de prueba separado.

\subsubsection{Liga de \textbf{GitHub}}

\href{https://github.com/MarcosHCK/master-programming/tree/master/clase011}{Aquí \faGithubSquare}

\subsection{Programa 3: \texttt{logicgates.ino}}

\subsubsection{Objetivos}

Implementar un perceptrón de dos entradas en Arduino UNO para reconocer la
función lógica AND. El resultado se muestra en un LED y en el Monitor Serie.

\subsubsection{Tecnicismos}

Arquitectura basada en \ref{sec:genarc:arduino}. Las entradas binarias son
\(X_1\) y \(X_2\), con el sesgo representado por \(X_0=1\). El perceptrón
calcula \(Z=W_0X_0+W_1X_1+W_2X_2\) y produce uno cuando \(Z\geq 0\). Los
pulsadores están conectados a D8 y D7 con \texttt{INPUT\_PULLUP}, por lo que
al pulsarlos la lectura se invierte a uno. La salida controla el LED en D4.

\subsubsection{Código relevante}

El sketch completo configura los pines, lee los pulsadores y actualiza la salida
del perceptrón:

\inputminted{cpp}{../clase011/logicgates.ino}

\subsubsection{Ejecución}

Cargar el sketch a un Arduino UNO R3. Conectar un pulsador entre D8 y GND y otro
entre D7 y GND; conectar el LED a D4 con su resistencia limitadora. Abrir el
Monitor Serie a \(9600\) baudios para ver los pesos y los cambios de entrada y
salida.

\subsubsection{Resultados esperados}

El LED se enciende únicamente cuando ambos pulsadores están presionados. La
tabla resume la tabla de verdad:

\begin{center}
\begin{tabular}{cc|c}
  \(X_1\) & \(X_2\) & \(Y\) \\
  \hline
  0 & 0 & 0 \\
  0 & 1 & 0 \\
  1 & 0 & 0 \\
  1 & 1 & 1 \\
\end{tabular}
\end{center}

\subsubsection{Liga de \textbf{GitHub}}

\href{https://github.com/MarcosHCK/master-programming/tree/master/clase011}{Aquí \faGithubSquare}

\subsection{Programa 4: \texttt{diabetes.ino}}

\subsubsection{Objetivos}

Aplicar en Arduino UNO el perceptrón entrenado con \texttt{diabetes.csv} para
clasificar registros según el índice de masa corporal (IMC) y la glucosa.

\subsubsection{Tecnicismos}

El modelo tiene dos entradas y usa el sesgo y los pesos reportados por
\texttt{clase011}. Cada entrada se normaliza con la media \(\mu_i\) y la escala
\(s_i\) del CSV; luego se calcula \(Z=W_0+\sum_i W_i((x_i-\mu_i)/s_i)\). La
salida es uno si \(Z\geq0\). El programa valida que ambos valores sean números
positivos.

\subsubsection{Código relevante}

El sketch completo contiene los parámetros aprendidos, el análisis de la línea
serie y la inferencia:

\inputminted{cpp}{../clase011/diabetes.ino}

\subsubsection{Ejecución}

Cargar el sketch a Arduino UNO y abrir el Monitor Serie a \(9600\) baudios.
Enviar una línea con el formato \texttt{IMC,Glucosa}, por ejemplo
\texttt{22.49,95.12}. El programa imprime la suma ponderada y la clase, y queda
listo para recibir otro registro.

\subsubsection{Resultados esperados}

Para el ejemplo \texttt{22.49,95.12}, el modelo produce la clase cero (sin
diabetes en las etiquetas del conjunto de datos). La etiqueta es una salida del
modelo académico y no una evaluación clínica.

\subsubsection{Liga de \textbf{GitHub}}

\href{https://github.com/MarcosHCK/master-programming/tree/master/clase011}{Aquí \faGithubSquare}

\subsection{Programa 5: \texttt{anemia.ino}}

\subsubsection{Objetivos}

Ejecutar en Arduino UNO el perceptrón de cuatro entradas entrenado con
\texttt{anemia.csv}. El sketch ilustra cómo llevar a una placa un clasificador
entrenado fuera de ella; su salida tiene propósito académico y no diagnostica
anemia.

\subsubsection{Tecnicismos}

Las entradas son edad, energía, frecuencia y antecedentes. El firmware contiene
el sesgo, los cuatro pesos y las medias y escalas del conjunto de entrenamiento.
Normaliza cada característica y calcula una suma ponderada; la clase es uno
cuando la suma es no negativa. Antecedentes debe ser cero o uno, la edad y la
frecuencia deben ser positivas y la energía no puede ser negativa.

\subsubsection{Código relevante}

El listado incluye los pesos aprendidos, las constantes de normalización, la
validación de la línea CSV y la clasificación:

\inputminted{cpp}{../clase011/anemia.ino}

\subsubsection{Ejecución}

Cargar el sketch a Arduino UNO y abrir el Monitor Serie a \(9600\) baudios.
Introducir cuatro valores separados por comas, en este orden: edad, energía,
frecuencia y antecedentes. Por ejemplo: \texttt{25,8,72,0}. El sketch imprime
\(Z\) y la clase predicha.

\subsubsection{Resultados esperados}

El registro de ejemplo pertenece a la clase cero (sin anemia en los datos). Una
entrada con antecedentes debe terminar en cero o uno; el resultado es una
clasificación del modelo, no un diagnóstico.

\subsubsection{Liga de \textbf{GitHub}}

\href{https://github.com/MarcosHCK/master-programming/tree/master/clase011}{Aquí \faGithubSquare}

\subsection{Programa 6: \texttt{obesidad.ino}}

\subsubsection{Objetivos}

Aplicar en Arduino UNO el perceptrón de seis entradas entrenado con
\texttt{obesidad.csv}, usando datos introducidos por el Monitor Serie.

\subsubsection{Tecnicismos}

El modelo usa edad, sexo, IMC, actividad, calorías y antecedentes. El sketch
incluye el sesgo, los pesos y las medias y escalas aprendidos; normaliza las seis
entradas antes de calcular \(Z\), y asigna la clase uno cuando \(Z\geq0\). Sexo
y antecedentes deben ser binarios; actividad debe ser no negativa y edad, IMC y
calorías deben ser positivos.

\subsubsection{Código relevante}

El sketch completo implementa la validación de los seis campos, la normalización
y la inferencia:

\inputminted{cpp}{../clase011/obesidad.ino}

\subsubsection{Ejecución}

Cargar el sketch a Arduino UNO y abrir el Monitor Serie a \(9600\) baudios.
Enviar los valores en el orden
edad, sexo, IMC, actividad, calorías y antecedentes; por ejemplo,
\texttt{25,0,22.5,5,2200,0}. El programa imprime la suma ponderada y la clase.

\subsubsection{Resultados esperados}

Para el registro de ejemplo, el modelo predice la clase cero (sin obesidad en las
etiquetas del CSV).

\subsubsection{Liga de \textbf{GitHub}}

\href{https://github.com/MarcosHCK/master-programming/tree/master/clase011}{Aquí \faGithubSquare}

\end{document}
